% Part: turing-machines
% Chapter: undecidability
% Section: verification

\documentclass[../../../include/open-logic-section]{subfiles}

\begin{document}

\olfileid{tur}{und}{ver}
\olsection{નિરૂપણની ચકાસણી}

\begin{explain}
આપણું નિરૂપણ કામ કરે છે તેની ચકાસણી માટે બે વાત
સાબિત કરવી પડશે. પ્રથમ, $M$ ઇનપુટ~$w$ પર અટકે,
તો $!T(M, w) \lif !E(M, w)$ પ્રામાણ્ય ધરાવે છે તે
બતાવવું પડશે. પછી ઊલટી દિશા: જો
$!T(M, w) \lif !E(M, w)$ પ્રામાણ્ય ધરાવે, તો
ઇનપુટ~$w$ પર ચલાવતાં $M$ ખરેખર છેવટે અટકે છે.

આ બે પરિણામો સાબિત કરવાની રીતો એકદમ જુદી છે.
પહેલા પરિણામ માટે પ્રથમ-ક્રમના તર્કનું !!a{sentence}
(એટલે $!T(M, w) \lif !E(M, w)$) પ્રામાણ્ય ધરાવે છે
તે બતાવવું છે. તેનો સૌથી સરળ રસ્તો !!a{derivation}
આપવાનો છે. પરંતુ આપણી સાબિતી બધાં~$M$ અને~$w$
માટે કામ કરવી જોઈએ; એટલે એક જ !!{sentence}ની
!!a{derivation} આપવાની નથી, અસંખ્ય વાક્યો માટે
આપવાની છે. તેથી $M$ અને~$w$ ગમે તે હોય અને
ઇનપુટ~$w$ પર અટકતાં~$M$ને ગમે તેટલાં પગલાં લાગતાં
હોય, $!T(M, w) \lif !E(M, w)$ની !!a{derivation}
મળશે એવું આગમનથી સાબિત કરી શકીએ.

સ્વાભાવિક રીતે આગમનનો આધાર~$M$ અટકવાના સંરૂપણ
સુધી પહોંચે તે પહેલાંના પગલાંની સંખ્યા હશે. આ
સાબિતીમાં આપણે સ્થાપીશું કે ઇનપુટ~$w$ પરના~$M$ના
સંગણનના દરેક પગલા~$n$ માટે
$!T(M, w) \Entails !C(M, w, n)$, જ્યાં
$!C(M, w, n)$ એ ઇનપુટ~$w$ પર ચાલતા~$M$નું
$n$ પગલાં પછીનું સંરૂપણ યોગ્ય રીતે વર્ણવે છે.
હવે માનો કે $M$ ઇનપુટ~$w$
પર~$n$ પગલાં પછી અટકે છે. ત્યારે
$!C(M, w, n)$ અટકવાનું સંરૂપણ વર્ણવશે. જ્યારે
$!C(M, w, n)$ અટકવાનું સંરૂપણ વર્ણવે, ત્યારે
$!C(M, w, n) \Entails !E(M, w)$ પણ બતાવીશું.
  એટલે $M$ ઇનપુટ~$w$ પર અટકે, તો કોઈ~$n$ માટે
$n$ પગલાં પછી~$M$ અટકવાના સંરૂપણમાં હશે. તેથી
$!T(M, w) \Entails !C(M, w, n)$, જ્યાં
$!C(M, w, n)$ એવું સંરૂપણ વર્ણવે છે; અને એ કિસ્સામાં
$!C(M, w, n) \Entails !E(M, w)$ હોવાથી
$!T(M, w) \Entails !E(M, w)$ મળે છે, એટલે
$\Entails !T(M, w) \lif !E(M, w)$.
\sourcecorrection{OLUN-006}{સ્થિર અંગ્રેજી મૂળની આ
છેલ્લી ફલિતતા-પંક્તિમાં વ્યાખ્યાયિત $!T(M,w)$ને
બદલે $T(M,w)$ છપાયું છે; અહીં વ્યાખ્યાયિત
વાક્યનું ચિહ્ન જાળવ્યું છે.}

ઊલટી દિશાની રીત તદ્દન જુદી છે. અહીં આપણે
$\Entails !T(M, w) \lif !E(M, w)$ માનીને
$M$ ઇનપુટ~$w$ પર અટકે છે તે સાબિત કરવાનું છે.
ધારણામાંથી $!T(M, w) \Entails !E(M, w)$ મળે છે:
એટલે જ્યાં $!T(M, w)$ સત્ય છે એવી દરેક
!!{structure}માં $!E(M, w)$ પણ સત્ય છે. તેથી
જેમાં $!T(M, w)$ સત્ય છે એવી
!!a{structure}~$\Struct{M}$ વર્ણવીશું. તેનો પ્રદેશ
$\Nat$ હશે; બધાં $\Obj Q_q$ અને $\Obj S_\sigma$
નું અર્થઘટન ઇનપુટ~$w$ પર ચાલતા~$M$નાં
સંરૂપણો પરથી નક્કી થશે. દાખલા તરીકે,
$\Sat{M}{\Obj Q_q(\num{m}, \num{n})}$ ત્યારે અને માત્ર
ત્યારે જ્યારે~$M$ ઇનપુટ~$w$ પર~$n$ પગલાં ચાલ્યા
પછી અવસ્થા~$q$માં હોય અને ખાનું~$m$ વાંચતું હોય.
\sourcecorrection{OLUN-007}{સ્થિર અંગ્રેજી મૂળના આ
દાખલામાં ચાલતા મશીનને $T$ લખ્યું છે; સંદર્ભમાં તે
$M$ છે, જ્યારે $!T(M,w)$ મશીનનું વર્ણન કરતું વાક્ય છે.}
હવે ધારણા મુજબ $!T(M, w) \Entails !E(M, w)$,
અને રચના મુજબ $\Sat{M}{!T(M, w)}$;
તેથી $\Sat{M}{!E(M, w)}$. પરંતુ
$\Sat{M}{!E(M, w)}$ ત્યારે અને માત્ર ત્યારે જ્યારે
કોઈ $n \in \Domain{M} = \Nat$ માટે ઇનપુટ~$w$
પર ચાલતું~$M$,~$n$ પગલાં પછી અટકવાના
સંરૂપણમાં હોય.
\end{explain}

\begin{defn}
$!C(M, w, n)$ એ નીચેનું !!{sentence} માનો:
\[
\Obj Q_q(\num{m}, \num{n}) \land \Obj S_{\sigma_0}(\num{0}, \num{n})
\land \dots \land \Obj S_{\sigma_k}(\num{k}, \num{n}) \land
\lforall[x][(\num{k} < x \lif \Obj S_\TMblank(x, \num{n}))]
\]
અહીં~$q$, સમય~$n$ પર~$M$ની અવસ્થા છે;
સમય~$n$ પર~$M$ ખાનું~$m$ વાંચે છે;
$0 \le i \le k$ માટે ખાનું~$i$ સમય~$n$ પર
ચિહ્ન~$\sigma_i$ ધરાવે છે;
અને~$k$ એ સમય~$0$ પરની ટેપનું સૌથી જમણું ખાલી
ન હોય એવું ખાનું તથા~$n$ પગલાં પછી હેડે મુલાકાત
લીધેલું સૌથી જમણું ખાનું—આ બેમાંથી વધુ જમણાં
ખાનાનો સૂચક છે.
\end{defn}

\begin{lem}\ollabel{lem:halt-config-implies-halt}
ઇનપુટ~$w$ પર ચાલતું~$M$,~$n$ પગલાં પછી
અટકવાના સંરૂપણમાં હોય, તો
$!C(M, w, n) \Entails !E(M, w)$.
\end{lem}

\begin{proof}
માનો કે ઇનપુટ~$w$ પર~$M$,~$n$ પગલાં પછી અટકે છે.
કોઈ અવસ્થા~$q$, ખાનું~$m$ અને ચિહ્ન~$\sigma$
એવાં છે કે:
\begin{enumerate}
\item $n$ પગલાં પછી~$M$ અવસ્થા~$q$માં ખાનું~$m$
  વાંચે છે, અને તે ખાનામાં~$\sigma$ છે.
\item સંક્રમણ વિધેય~$\delta(q, \sigma)$ અવ્યાખ્યાયિત છે.
\end{enumerate}
$!C(M, w, n)$ આ સંરૂપણનું વર્ણન છે અને તેમાં
$\Obj Q_{q}(\num{m}, \num{n})$ તથા
$\Obj S_{\sigma}(\num{m}, \num{n})$ ખંડો હશે.
આ બંને ખંડો મળીને~$!E(M, w)$ ફલિત કરે છે:
\[
\lexists[x][\lexists[y][(\bigvee_{\tuple{q, \sigma} \in
      X}(\Obj Q_q(x, y) \land \Obj S_\sigma(x, y)))]]
\]
કારણ કે $\tuple{q,\sigma} \in X$ હોવાથી
$\Obj Q_q(\num{m}, \num{n}) \land
\Obj S_{\sigma}(\num{m}, \num{n}) \Entails
\bigvee_{\tuple{q, \sigma} \in X}
(\Obj Q_q(\num{m}, \num{n}) \land
\Obj S_{\sigma}(\num{m}, \num{n}))$.
\sourcecorrection{OLUN-008}{સ્થિર અંગ્રેજી મૂળની
આ ફલિતતા-પંક્તિમાં પહેલાના પસંદ કરેલા $q,\sigma$
બદલે અચાનક $q',\sigma'$ આવે છે અને એક
$S$-વિધેય આગળનો વસ્તુભાષાનો અક્ષરપ્રકાર આપતો આદેશ છૂટી ગયો છે. અહીં
અગાઉના જ $q,\sigma$ અને પૂરો વિધેય-સંકેત રાખ્યો છે.}
\end{proof}

\begin{explain}
એટલે $M$ ઇનપુટ~$w$ પર અટકે, તો કોઈ~$n$ માટે
$!C(M, w, n) \Entails !E(M,w)$. હવે આપણે
બતાવીશું કે કોઈ પણ સમય~$n$ માટે
$!T(M, w) \Entails !C(M, w, n)$.
\end{explain}

\begin{lem}
\ollabel{lem:config}
દરેક~$n$ માટે, $M$ જો~$n$ પગલાં પૂરાં થાય તે
પહેલાં અટક્યું ન હોય, તો
$!T(M, w) \Entails !C(M, w, n)$.
\sourcecorrection{OLUN-009}{સ્થિર અંગ્રેજી મૂળનું
``$n$ પગલાં પછી અટક્યું ન હોય'' વિધાન~$n$મા
પગલે જ અટકવાનું સંરૂપણ બાકાત કરે છે, છતાં નીચેનો
પ્રામાણ્યનો પુરાવો એ જ સંરૂપણ માટે આ ઉપપ્રમેય
વાપરે છે. અહીં~$n$ પગલાં પૂરાં થતાં પહેલાં ન
અટકવાની જરૂરી શરત સ્પષ્ટ કરી છે.}
\end{lem}

\begin{proof}
આગમનનો આધાર: $n = 0$ હોય, તો
$!C(M, w, 0)$ના સંયોજકો~$!T(M,w)$નાં
પ્રારંભિક સંરૂપણ અને સંખ્યાનાં સ્વયંસિદ્ધ વાક્યોમાંથી
ફલિત થાય છે. ખાલી ઇનપુટમાં શરૂઆતથી વાંચાતું
ખાનું ખાલી છે, અને તેની જમણી બાજુનાં બધાં
ખાનાં પણ ખાલી છે, તે ખાલી-ટેપની શરત અને
અનુગામીના સ્વયંસિદ્ધ વાક્યથી મળે છે.
\sourcecorrection{OLUN-024}{સ્થિર અંગ્રેજી મૂળ
આધારકિસ્સાના બધા સંયોજકો~$!T(M,w)$માં
સીધા જ છે એમ કહે છે. ખાલી ઇનપુટ હોય ત્યારે
માથું શરૂઆતમાં પહેલું ખાનું વાંચે છે, એટલે
તે ખાનું ખાલી હોવાનું અને તેની આગળની
ખાલી-ટેપ શરતનું અનુમાન પણ કરવું પડે છે.}

આગમનની ધારણા: $M$ જો~$n$મા પગલા પહેલાં
અટક્યું ન હોય, તો $!T(M,w) \Entails !C(M, w, n)$.
હવે બતાવવું છે કે, $!C(M, w, n)$ અટકવાનું સંરૂપણ
વર્ણવતું ન હોય તો,
$!T(M, w) \Entails !C(M, w, n+1)$.

માનો કે $n \ge 0$ અને ઇનપુટ~$w$ પર શરૂ થયેલું
$M$,~$n$ પગલાં પછી અવસ્થા~$q$માં ખાનું~$m$
વાંચે છે. $M$,~$n$ પગલાં પછી અટકતું ન હોવાથી
~$M$ના કાર્યક્રમમાં નીચેના ત્રણ સ્વરૂપોમાંની કોઈ એક
સૂચના હોવી જોઈએ:
\sourcecorrection{OLUN-010}{સ્થિર અંગ્રેજી મૂળ આ
આગમન-પગલામાં $n>0$ ધારે છે, જેથી આધાર પછીનું
$n=0$થી $n+1$ પગલું છૂટી જાય છે. એ જ વિચાર
$n=0$ માટે પણ લાગુ પડે છે, તેથી અહીં $n\ge0$ છે.}

\begin{enumerate}
\item \ollabel{right} $\delta(q, \sigma) = \tuple{q', \sigma', \TMright}$

\item \ollabel{left} $\delta(q, \sigma) = \tuple{q', \sigma', \TMleft}$

\item \ollabel{stay} $\delta(q, \sigma) = \tuple{q', \sigma', \TMstay}$
\end{enumerate}

આ ત્રણેય કિસ્સા અનુક્રમે જોઈએ.

\begin{enumerate}
\item માનો કે સૂચના~\olref{right}ના સ્વરૂપની છે.
  \olref[rep]{defn:tm-descr}\olref[rep]{rep-right}
  મુજબ આનો અર્થ છે કે
\begin{align*}
& \lforall[x][\lforall[y][((\Obj Q_{q}(x,
  y) \land \Obj S_\sigma(x, y)) \lif {}]]\\
& \qquad (\Obj Q_{q'}(x',
  y') \land \Obj S_{\sigma'}(x, y') \land !A(x, y)))
\intertext{એ $!T(M,w)$નો એક સંયોજક છે. તેથી
  સર્વપરિમાણકના ઉદાહરણથી (અહીં~$x$ માટે
  $\num{m}$ અને~$y$ માટે $\num{n}$ મૂકીને)
  નીચેનું !!{sentence} ફલિત થાય છે:}
& (\Obj Q_{q}(\num{m}, \num{n}) \land \Obj S_{\sigma}(\num{m},
\num{n})) \lif {}\\
& \qquad (\Obj Q_{q'}(\num{m}', \num{n}') \land
\Obj S_{\sigma'}(\num{m}, \num{n}') \land !A(\num{m}, \num{n})).
\intertext{આગમનની ધારણા મુજબ $!T(M, w) \Entails !C(M, w, n)$,
  એટલે કે:}
& \Obj Q_q(\num{m}, \num{n}) \land \Obj S_{\sigma_0}(\num{0}, \num{n})
\land \dots \land \Obj S_{\sigma_k}(\num{k}, \num{n}) \land \\
& \qquad\lforall[x][(\num{k} < x \lif \Obj S_\TMblank(x, \num{n}))]\\
\intertext{$n$ પગલાં પછી ખાનું~$m$ ચિહ્ન~$\sigma$
  ધરાવે છે, તેથી અનુરૂપ સંયોજક
  $\Obj S_\sigma(\num{m}, \num{n})$ છે; એટલે
  તેમાંથી આ ફલિત થાય છે:}
& \Obj Q_{q}(\num{m}, \num{n}) \land \Obj S_{\sigma}(\num{m},
   \num{n})
\intertext{હવે આપણને મળે છે:}
& \Obj Q_{q'}(\num{m}', \num{n}') \land \Obj S_{\sigma'}(\num{m},
  \num{n}') \land {}\\
& \qquad\bigwedge_{\substack{0\le i\le k\\i\ne m}}
  \Obj S_{\sigma_i}(\num{i}, \num{n}') \land {}\\
& \qquad \lforall[x][(\num{k} < x \lif \Obj S_\TMblank(x, \num{n}'))]
\end{align*}
\sourcecorrection{OLUN-022}{સ્થિર અંગ્રેજી મૂળના
જમણી-ચાલના પરિણામમાં બદલાયેલા ખાના~$m$નું
નવું ચિહ્ન પહેલેથી અલગ આપેલું છે, છતાં પછીની
યાદી દેખીતી રીતે $0$થી~$k$ સુધીના બધા
જૂનાં ચિહ્નો ફરી જોડે છે. મૂળનું જ ગદ્ય~$m$નો
બાકાત રાખે છે; અહીં એ બાકાત સ્પષ્ટ લખ્યો છે.}
આ રીતે: પહેલી પંક્તિ અગાઉના શરતી વાક્યના
અનુગામીમાંથી સીધી, મોડસ પોનેન્સથી મળે છે.
વચ્ચેની પંક્તિનો દરેક સંયોજક—જેમાં
$S_{\sigma_m}(\num{m},\num{n}')$નો સમાવેશ નથી—
$!C(M, w, n)$ના અનુરૂપ સંયોજક તથા
$!A(\num{m}, \num{n})$માંથી ફલિત થાય છે.

$m < k$ હોય, તો $!T(M,w) \Proves \num{m} < \num {k}$
(\olref[rep]{prop:mlessk}); અને~$<$ની સંક્રામકતાથી
$\lforall[x][(\num{k} < x \lif \num{m} < x)]$
મળે છે. $m = k$ હોય, તો એ જ વાક્ય
$\lforall[x][(\num{k} < x \lif \num{m} < x)]$
માત્ર તર્કથી મળે છે. તેથી છેલ્લી પંક્તિ
$!C(M, w, n)$ના અનુરૂપ સંયોજક,
$\lforall[x][(\num{k} < x \lif \num{m} < x)]$,
અને $!A(\num{m}, \num{n})$માંથી ફલિત થાય છે.
$m<k$ હોય, તો આ જ $!C(M, w, n+1)$ છે.

હવે માનો કે $m=k$. ત્યારે $n+1$ પગલાં પછી
ટેપ-હેડે ખાનું~$k+1$ પણ મુલાકાત લીધું હશે;
તે હવે મુલાકાત લેવાયેલું સૌથી જમણું ખાનું છે.
એટલે $!C(M, w, n+1)$માં નવો સંયોજક
$\Obj S_\TMblank(\num{k}',\num{n}')$ છે, અને
છેલ્લો સંયોજક
$\lforall[x][(\num{k}' < x \lif \Obj S_\TMblank(x, \num{n}'))]$
છે. આ બે !!{sentence}s પણ ફલિત થાય છે તે
ચકાસવું પડશે.

આપણી પાસે પહેલેથી
$\lforall[x][(\num{k} < x \lif \Obj S_\TMblank(x,
  \num{n}'))]$ છે. ખાસ કરીને તેમાંથી
$\num{k} < \num{k}' \lif
\Obj S_\TMblank(\num{k}', \num{n}')$ મળે છે.
સ્વયંસિદ્ધ વાક્ય $\lforall[x][x < x']$માંથી
$\num{k} < \num{k}'$ મળે છે. મોડસ પોનેન્સથી
$\Obj S_\TMblank(\num{k}',\num{n}')$ ફલિત થાય છે.

વળી, $!T(M,w) \Proves \num{k} < \num{k}'$
હોવાથી~$<$ની સંક્રામકતાના સ્વયંસિદ્ધ વાક્યથી
$\lforall[x][(\num{k}' < x \lif \Obj
  S_\TMblank(x, \num{n}'))]$ મળે છે. (આ
ચકાસણી કસરત તરીકે છોડી છે.)

\item માનો કે સૂચના~\olref{left}ના સ્વરૂપની છે.
  તો \olref[rep]{defn:tm-descr}\olref[rep]{rep-left}
  મુજબ:
\sourcecorrection{OLUN-011}{સ્થિર અંગ્રેજી મૂળના
આ ડાબે-ચાલ કિસ્સામાં લખાતું ખાનું~$x'$ અને
ઉદાહરણમાં~$l+1$ છે, પરંતુ યથાવત્-ખાનાંની
શરત અનુક્રમે $!A(x,y)$ અને $!A(\num{l},\num{n})$
લખાઈ છે. ઉપરાંત શૂન્ય-ખાનાની શાખામાં અહીં
પસંદ કરેલા $q,q'$ને બદલે જૂનાં $q_i,q_j$ છે.
નીચે લખાતું ખાનું બાકાત રાખ્યું છે અને
અવસ્થા-ચિહ્નો એકરૂપ કર્યા છે.}
\begin{align*}
& \lforall[x][\lforall[y][((\Obj Q_{q}(x', y) \land \Obj
    S_{\sigma}(x', y)) \lif {}]]\\
& \qquad (\Obj Q_{q'}(x, y') \land \Obj
  S_{\sigma'}(x', y') \land !A(x', y))) \land {}\\
& \lforall[y][((\Obj Q_{q}(\num{0}, y) \land \Obj S_{\sigma}(\num{0},
    y)) \lif {}]\\
& \qquad (\Obj Q_{q'}(\num{0}, y') \land \Obj S_{\sigma'}(\num{0}, y')
  \land !A(\num{0}, y)))
\intertext{એ $!T(M,w)$નો એક સંયોજક છે. $m>0$
  હોય, તો $l = m - 1$ લો (એટલે $m = l+1$).
  ઉપરના !!{sentence}ના પહેલા સંયોજકમાંથી
  નીચેનું ફલિત થાય છે:}
& (\Obj Q_{q}(\num{l}', \num{n}) \land \Obj S_{\sigma}(\num{l}', \num{n}))
\lif {} \\
& \qquad (\Obj Q_{q'}(\num{l}, \num{n}') \land \Obj S_{\sigma'}(\num{l}',
\num{n}') \land !A(\num{l}', \num{n}))
\intertext{અન્યથા $l = m = 0$ લો અને બીજા
  સંયોજકમાંથી ફલિત થતું નીચેનું !!{sentence}
  જુઓ:}
& ((\Obj Q_{q}(\num{0}, \num{n}) \land \Obj S_{\sigma}(\num{0}, \num{n})) \lif {}\\
& \qquad   (\Obj Q_{q'}(\num{0}, \num{n}') \land \Obj S_{\sigma'}(\num{0}, \num{n}') \land
!A(\num{0}, \num{n})))
\intertext{બંનેમાંથી જે લાગુ પડે તેમાંથી
નીચેનું ફલિત થાય છે:}
&  \Obj Q_{q'}(\num{l}, \num{n}') \land \Obj S_{\sigma'}(\num{m},
  \num{n}') \land {}\\
&\qquad  \bigwedge_{\substack{0\le i\le k\\i\ne m}}
  \Obj S_{\sigma_i}(\num{i}, \num{n}') \land {} \\
&\qquad  \lforall[x][(\num{k} < x
  \lif \Obj S_\TMblank(x, \num{n}'))]
\end{align*}
\sourcecorrection{OLUN-023}{સ્થિર અંગ્રેજી મૂળના
ડાબી-ચાલના પરિણામમાં પણ લખાયેલું ખાનું~$m$
અલગ નવા ચિહ્ન સાથે આવે છે; પછીની જૂનાં
ચિહ્નોની યાદીમાં તેને ફરી સામેલ કરવાથી
સંરૂપણ ખોટું થાય. અહીં તે ખાનું સ્પષ્ટ બાકાત છે.}
અગાઉની જેમ. (પહેલા કિસ્સામાં
$\num{l}' \ident \num{l+1} \ident \num{m}$
અને બીજા કિસ્સામાં $\num{l} \ident \num{0}$
એ નોંધો.) પરંતુ આ તો $!C(M, w, n+1)$ જ છે.

\item \olref{stay} કિસ્સો કસરત તરીકે છોડી છે.
\end{enumerate}
આપણે બતાવ્યું કે આવા દરેક~$n$ માટે
$!T(M, w) \Entails !C(M, w, n)$.
\end{proof}

\begin{prob}
\olref[tur][und][ver]{lem:config}ની સાબિતીમાં
કિસ્સો~\olref[tur][und][ver]{stay} પૂર્ણ કરો.
\end{prob}

\begin{prob}
$\Obj S_{\sigma_i}(\num{i}, \num{n})$ અને
$!A(m, n)$માંથી
$\Obj S_{\sigma_i}(\num{i}, \num{n}')$ની
!!a{derivation} આપો (ધારો કે $i \neq m$,
એટલે $i < m$ અથવા $m < i$).
\end{prob}

\begin{prob}
$\lforall[x][(\num{k} < x \lif \Obj
  S_\TMblank(x, \num{n}'))]$,
$\lforall[x][x < x']$ અને
$\lforall[x][\lforall[y][\lforall[z][ ((x < y \land y < z) \lif x <
      z)]]]$માંથી
$\lforall[x][(\num{k}' < x \lif \Obj
  S_\TMblank(x, \num{n}'))]$ની
!!a{derivation} આપો.
\end{prob}

\begin{lem}
\ollabel{lem:valid-if-halt}
$M$ ઇનપુટ~$w$ પર અટકે, તો
$!T(M, w) \lif !E(M, w)$ પ્રામાણ્ય ધરાવે છે.
\end{lem}

\begin{proof}
\olref{lem:config}થી જાણીએ છીએ કે પ્રથમ
અટકવાના સમય સુધીનાં દરેક~$n$ માટે
સમય~$n$ પરનાં~$M$ના સંરૂપણનું
વર્ણન~$!C(M, w, n)$,
$!T(M, w)$માંથી ફલિત થાય છે.
\sourcecorrection{OLUN-012}{સ્થિર અંગ્રેજી મૂળ
અહીં ``કોઈ પણ સમય~$n$'' કહે છે, છતાં મશીન
અટક્યા પછીનું સંરૂપણ તેની વ્યાખ્યા મુજબ આપેલું
નથી. OLUN-009ના સુધારેલા ઉપપ્રમેયને પ્રથમ
અટકવાના સમય સુધી જ વાપર્યો છે.}
માનો કે $M$,~$k$ પગલાં પછી અટકે છે. તે સમયે
તે કોઈ~$m \in \Nat$ માટે ખાનું~$m$ વાંચતું હશે.
તેથી $!C(M, w, k)$ એ~$M$ના અટકવાના
સંરૂપણનું વર્ણન છે: તેમાં
$\Obj Q_q(\num{m}, \num{k})$ તથા
$\Obj S_\sigma(\num{m}, \num{k})$ બંને સંયોજકો છે,
અને $\delta(q,\sigma)$ અવ્યાખ્યાયિત છે. તેથી
\olref{lem:halt-config-implies-halt} મુજબ
$!C(M, w, k) \Entails !E(M, w)$. વળી
$!T(M, w) \Entails !C(M, w, k)$ હોવાથી
$!T(M, w) \Entails !E(M, w)$; એટલે
$!T(M, w) \lif !E(M, w)$ પ્રામાણ્ય ધરાવે છે.
\end{proof}

\begin{explain}
આપણો દાવો ચકાસવાનું પૂરું કરવા ઊલટી દિશા પણ
સ્થાપવી પડશે: $!T(M, w) \lif !E(M, w)$
પ્રામાણ્ય ધરાવે, તો ઇનપુટ~$w$ પર શરૂ કરતાં
$M$ ખરેખર અટકે છે.
\end{explain}

\begin{lem}
\ollabel{lem:halt-if-valid}
$\Entails !T(M, w) \lif !E(M, w)$ હોય, તો
$M$ ઇનપુટ~$w$ પર અટકે છે.
\end{lem}

\begin{proof}
પ્રદેશ~$\Nat$ ધરાવતું $\Lang L_M$-!!{structure}
$\Struct M$ લો. તેમાં $\Obj 0$નું અર્થઘટન~$0$,
$\prime$નું અનુગામી વિધેય અને~$<$નું નાની-કરતાં
સંબંધ તરીકે કરીએ છીએ. વિધાનો $\Obj Q_q$ અને
$\Obj S_\sigma$નું અર્થઘટન આ પ્રમાણે છે:
\begin{align*}
  \Assign{\Obj Q_q}{M} & =
\Setabs{\tuple{m, n}}{\begin{array}{ll}\text{ઇનપુટ~$w$ પર શરૂ કરતાં,~$n$ પગલાં પછી,}\\ \text{$M$ અવસ્થા~$q$માં
  ખાનું~$m$ વાંચે છે}\end{array}} \\
  \Assign{\Obj S_\sigma}{M} & = \Setabs{\tuple{m, n}}{\begin{array}{ll}
\text{ઇનપુટ~$w$ પર શરૂ કરતાં,~$n$ પગલાં પછી,}\\ \text{$M$ના ખાના~$m$માં
  ચિહ્ન~$\sigma$ છે}\end{array}}
\end{align*}
બીજા શબ્દોમાં, ઇનપુટ~$w$ પર શરૂ થયેલું~$M$
હકીકતમાં પગલે પગલે શું કરે છે તેનું વર્ણન કરે
એવું !!{structure}~$\Struct{M}$ રચીએ છીએ. સ્પષ્ટપણે
$\Sat{M}{!T(M, w)}$. જો
$\Entails !T(M, w) \lif !E(M, w)$, તો
$\Sat{M}{!E(M, w)}$ પણ; એટલે:
\[
\Sat{M}{\lexists[x][\lexists[y][(\bigvee_{\tuple{q, \sigma} \in
      X}(\Obj Q_q(x, y) \land \Obj S_\sigma(x, y)))]]}.
\]
$\Domain{M} = \Nat$ હોવાથી કોઈ
$m$, $n \in \Nat$ તથા એવાં~$q$ અને~$\sigma$
છે કે $\Sat{M}{\Obj Q_q(\num{m}, \num{n}) \land
\Obj S_\sigma(\num{m}, \num{n})}$ અને
$\delta(q, \sigma)$ અવ્યાખ્યાયિત છે.
$\Struct M$ની વ્યાખ્યા મુજબ તેનો અર્થ એ છે કે
ઇનપુટ~$w$ પર શરૂ થયેલું~$M$,~$n$ પગલાં પછી
અવસ્થા~$q$માં છે અને ચિહ્ન~$\sigma$ વાંચે છે;
સંક્રમણ વિધેય અવ્યાખ્યાયિત છે, એટલે~$M$ અટક્યું છે.
\end{proof}

\end{document}
